% ============================================================================
% PART I: CONSTRUCTION
% Section 2: Mathematical Toolkit
% ============================================================================
%
% Purpose
%   Centralized math chapter so the reader can consult one place for every
%   structure used later. This follows the user's preference (Final QFT
%   Project Overview.md): "its better to have it centralized rather than
%   scattered bc im always kind of frustrated at having something mathy
%   defined and then not being able to see any consequences of that definition
%   until pages of physics."
%
% Scope (from tqft_observables_literature_review.md Sec. 1)
%   - Homotopy and large gauge transformations.
%   - Bundles, connections, curvature, characteristic classes.
%   - Holonomy, Wilson loops, flat connections.
%   - Symplectic reduction and moduli spaces.
%   - Bordisms and the Atiyah--Segal definition (stated here, used in Sec. 4).
%
% Status: EXTRACT from archive. Primary sources:
%   - archive_2026-05-01/topology_review.tex (main body)
%   - archive_2026-05-01/differential_forms_csimons_foundations_v2.tex (forms)
%   - writings/brian_chern_simons_theory.tex (moduli, holonomy rigor)
%
% Length target: ~15--20 pages
%
% Pedagogical references to cite
%   Nakahara (\cite{Nakahara2003GeometryTopologyPhysics}) -- bundles, char classes
%   Baez-Muniain (\cite{BaezMuniain1994GaugeFieldsKnotsGravity}) -- geometry+knots
%   Manton-Sutcliffe (\cite{MantonSutcliffeTopologicalSolitons}) -- soliton moduli
%   Weinberg (\cite{WeinbergClassicalSolutionsQFT}) -- classical solutions/moduli
%
% Structural references
%   Atiyah (\cite{Atiyah1988TQFT})                      -- bordism / axioms
%   Elitzur-Moore-Schwimmer-Seiberg (\cite{ElitzurMooreSchwimmerSeiberg1989CanonicalCS})
%                                                       -- symplectic quantization of CS
%   Witten (\cite{Witten1989Jones})                     -- knot/link context
%
% Subsection plan
%   2.1  Manifolds, forms, Hodge star, orientations, and conventions
%   2.2  Lie groups, Lie algebras, adjoint action, Killing form
%   2.3  Principal bundles and connections
%   2.4  Curvature, gauge transformations, and characteristic classes
%   2.5  Homotopy groups, winding, and $\pi_n(G)$
%   2.6  Holonomy, Wilson loops, and flat connections
%   2.7  Moduli spaces of flat connections and symplectic reduction
%   2.8  Bordism categories and the Atiyah--Segal axioms (statement)
%
% Writing notes
%   - Introduce standard definitions and equations before applying them
%     (project CLAUDE.md convention).
%   - Theorems and proofs short: statement, sketch, consequence.
%   - Each subsection ends with a brief physical reading so the chapter is
%     not a dry math dump.

\section{Mathematical Toolkit}
\label{sec:mathematical-toolkit}

% TODO: Migrate content from archive_2026-05-01/topology_review.tex and
%       differential_forms_csimons_foundations_v2.tex, supplemented with
%       brian_chern_simons_theory.tex for moduli / holonomy rigor.

\subsection{Forms, Manifolds, and Conventions}
\label{subsec:forms-conventions}

% TODO: differential forms, wedge, exterior derivative, Hodge star, orientation.
% Source: archive/differential_forms_csimons_foundations_v2.tex
% Cite: \cite{Nakahara2003GeometryTopologyPhysics}

\subsection{Lie Groups and Lie Algebras}
\label{subsec:lie-groups}

% TODO: G, g, adjoint action, Killing form, structure constants.
% Source: archive/topology_review.tex
% Cite: \cite{Nakahara2003GeometryTopologyPhysics}, \cite{BaezMuniain1994GaugeFieldsKnotsGravity}

\subsection{Principal Bundles, Connections, Curvature}
\label{subsec:bundles-connections}

% TODO: principal G-bundle, connection 1-form A, curvature F = dA + A wedge A.
% Characteristic classes (Chern, Pontryagin) as cohomology classes.
% Cite: \cite{Nakahara2003GeometryTopologyPhysics}

\subsection{Gauge Transformations and Characteristic Classes}
\label{subsec:gauge-char-classes}

% TODO: gauge transformations, Chern-Simons form as transgression,
% relation d omega_CS = tr(F wedge F).
% Cite: \cite{BaezMuniain1994GaugeFieldsKnotsGravity}

\subsection{Homotopy and Winding}
\label{subsec:homotopy-winding}

% TODO: pi_n(G), pi_3(SU(N)) = Z, pi_2(G/H) for monopole sectors.
% This machinery is used for level quantization in Sec. 3 and for monopoles in Sec. 7.
% Cite: \cite{Nakahara2003GeometryTopologyPhysics}

\subsection{Holonomy, Wilson Loops, and Flat Connections}
\label{subsec:holonomy-flat}

% TODO: holonomy P exp(integral A), flat condition F = 0,
% moduli M_flat(M,G) = Hom(pi_1(M),G)/G.
% Cite: \cite{Witten1989Jones}, \cite{BaezMuniain1994GaugeFieldsKnotsGravity}

\subsection{Moduli Spaces and Symplectic Reduction}
\label{subsec:moduli-symplectic}

% TODO: symplectic structure on M_flat(Sigma,G), geometric quantization setup.
% Cite: \cite{ElitzurMooreSchwimmerSeiberg1989CanonicalCS},
%       \cite{MantonSutcliffeTopologicalSolitons}, \cite{WeinbergClassicalSolutionsQFT}

\subsection{Bordism Categories and the Atiyah--Segal Axioms}
\label{subsec:bordism-atiyah-segal}

% TODO: Bord_n, symmetric monoidal functor Z: Bord_n -> Vect_k.
% State axioms here; CS theory realizes them in Sec. 4.
% Cite: \cite{Atiyah1988TQFT}, \cite{BirminghamBlauRakowskiThompson1991TopologicalFieldTheory}
